\chapter{Supplementary Materials}

This appendix provides supplementary information supporting the main analysis, including detailed variable definitions, additional robustness checks, extended mathematical derivations, and descriptions of the machine learning algorithms used in the study.

\section{Detailed Variable Definitions}

Table~\ref{tab:appendix_variables} provides detailed definitions and coding for all variables used in the analysis.\footnote{The positivity assumption \citep{westreich2010invited} requires that for every combination of covariates, there is a non-zero probability of both ITN use and non-use. We assess this empirically in Chapter 6. Missing data handling follows \citet{little2019statistical}.}

\begin{table}[htbp]
	\centering
	\caption{Detailed Variable Definitions and Coding}
	\label{tab:appendix_variables}
	\begin{tabular}{@{}llp{5cm}@{}}
		\toprule
		\textbf{Variable} & \textbf{Source} & \textbf{Coding} \\
		\midrule
		malaria & PR (hml32) & 1 = positive (RDT or microscopy), 0 = negative \\
		itn\_use & PR (hml12) & 1 = slept under ITN last night, 0 = did not use \\
		age\_months & PR (age\_years) & Age in months = age\_years $\times$ 12 (0-59 months) \\
		sex & PR (hv104) & Male (1), Female (2) \\
		wealth & HR (hv270) & 1=Poorest, 2=Poorer, 3=Middle, 4=Richer, 5=Richest \\
		residence & HR (hv025) & Urban (1), Rural (2) \\
		maternal\_ed & KR (v106) & 0=None, 1=Primary, 2=Secondary, 3=Higher \\
		rainfall & Geocov & Annual precipitation in millimeters \\
		temp\_mean & Geocov & Mean annual temperature in degrees Celsius \\
		evi & Geocov & Enhanced Vegetation Index (unitless) \\
		aridity & Geocov & Aridity index (higher = more arid) \\
		wet\_days & Geocov & Number of wet days per year ($\geq$ 1mm rainfall) \\
		population & Geocov & Population count from high-resolution population grids \\
		elevation & Geocov (2020 only) & Height above sea level in meters \\
		\bottomrule
	\end{tabular}
\end{table}

\section{Additional Descriptive Tables}

\subsection{Malaria Prevalence by Age and Sex}

Table~\ref{tab:appendix_age_sex} shows malaria prevalence stratified by age group and sex for the combined dataset.

\begin{table}[htbp]
	\centering
	\caption{Malaria Prevalence by Age Group and Sex (Combined Dataset)}
	\label{tab:appendix_age_sex}
	\begin{tabular}{@{}lccc@{}}
		\toprule
		\textbf{Age Group (months)} & \textbf{Male (\%)} & \textbf{Female (\%)} & \textbf{Total (\%)} \\
		\midrule
		0-11 & 2.8 & 2.3 & 2.6 \\
		12-23 & 5.1 & 4.3 & 4.7 \\
		24-35 & 3.9 & 3.4 & 3.7 \\
		36-47 & 5.5 & 5.0 & 5.3 \\
		48-59 & 4.2 & 3.9 & 4.1 \\
		\bottomrule
	\end{tabular}
\end{table}

\subsection{Environmental Variables by Survey Year}

Table~\ref{tab:appendix_env_by_year} presents environmental variables separately for each survey year.

\begin{table}[htbp]
	\centering
	\caption{Environmental Variables by Survey Year}
	\label{tab:appendix_env_by_year}
	\begin{tabular}{@{}lccc@{}}
		\toprule
		\textbf{Variable} & \textbf{KMIS 2015} & \textbf{KMIS 2020} & \textbf{Combined} \\
		\midrule
		Rainfall (mm) & 1,523 (698) & 1,589 (746) & 1,557 (723) \\
		Temperature ($^\circ$C) & 22.2 (2.8) & 22.0 (2.7) & 22.1 (2.8) \\
		EVI & 1,547 (1,721) & 1,598 (1,788) & 1,573 (1,755) \\
		Aridity & 29.5 (12.1) & 30.1 (12.5) & 29.8 (12.3) \\
		Wet days & 15.2 (4.1) & 15.0 (4.0) & 15.1 (4.1) \\
		\bottomrule
	\end{tabular}
	\caption*{Values are mean (standard deviation)}
\end{table}

\section{Additional Robustness Checks}

\subsection{Alternative Propensity Score Specification}

To assess sensitivity to the propensity score specification, we estimated PSM using different sets of covariates. Table~\ref{tab:appendix_psm_alt} shows the results.

\begin{table}[htbp]
	\centering
	\caption{PSM Results Under Alternative Covariate Specifications (Combined Dataset)}
	\label{tab:appendix_psm_alt}
	\begin{tabular}{@{}llccc@{}}
		\toprule
		\textbf{Specification} & \textbf{Covariates} & \textbf{OR} & \textbf{95\% CI} & \textbf{p-value} \\
		\midrule
		Minimal & Demographics only & 0.892 & (0.723-1.101) & 0.287 \\
		Medium & + Socioeconomic & 0.521 & (0.387-0.702) & $<$0.001 \\
		Full & + Environmental & 0.458 & (0.324-0.647) & $<$0.001 \\
		\bottomrule
	\end{tabular}
\end{table}

\subsection{Alternative Caliper for PSM}

We tested sensitivity to the matching caliper size. Following \citet{austin2010optimal}, we used a caliper of 0.2 standard deviations of the logit propensity score as the baseline. Table~\ref{tab:appendix_caliper} presents the results.

\begin{table}[htbp]
	\centering
	\caption{PSM Results Under Different Caliper Sizes (Combined Dataset)}
	\label{tab:appendix_caliper}
	\begin{tabular}{@{}lccc@{}}
		\toprule
		\textbf{Caliper (SD)} & \textbf{OR} & \textbf{95\% CI} & \textbf{Matched N} \\
		\midrule
		0.1 & 0.452 & (0.318-0.642) & 3,421 \\
		0.2 (baseline) & 0.458 & (0.324-0.647) & 3,980 \\
		0.3 & 0.461 & (0.329-0.646) & 4,102 \\
		No caliper & 0.465 & (0.335-0.645) & 4,215 \\
		\bottomrule
	\end{tabular}
\end{table}

\subsection{Alternative Number of Trees in Random Forest}

Table~\ref{tab:appendix_trees} shows DML results using different numbers of trees in the random forest learner.

\begin{table}[htbp]
	\centering
	\caption{DML Results with Different Numbers of Trees (Combined Dataset)}
	\label{tab:appendix_trees}
	\begin{tabular}{@{}lccc@{}}
		\toprule
		\textbf{Number of Trees} & \textbf{OR} & \textbf{95\% CI} & \textbf{p-value} \\
		\midrule
		100 & 0.978 & (0.959-0.997) & 0.023 \\
		250 & 0.980 & (0.962-0.998) & 0.031 \\
		500 (baseline) & 0.981 & (0.963-0.999) & 0.041 \\
		1000 & 0.979 & (0.961-0.998) & 0.035 \\
		\bottomrule
	\end{tabular}
\end{table}

\section{Extended DML Mathematical Details}

\subsection{The Partially Linear Model with Cross-Fitting}

Let $\{W_i = (Y_i, D_i, X_i)\}_{i=1}^n$ be independent and identically distributed observations. The partially linear model is:

\[
Y_i = \theta_0 D_i + g_0(X_i) + \varepsilon_i, \quad \mathbb{E}[\varepsilon_i \mid X_i, D_i] = 0
\]
\[
D_i = m_0(X_i) + \eta_i, \quad \mathbb{E}[\eta_i \mid X_i] = 0
\]

The nuisance functions $g_0(X) = \mathbb{E}[Y \mid X]$ and $m_0(X) = \mathbb{E}[D \mid X]$ are unknown and estimated using machine learning.

The DML estimator uses cross-fitting to avoid overfitting bias:

\begin{enumerate}
	\item Split the data into $K$ folds $\{I_k\}_{k=1}^K$ of approximately equal size.
	\item For each fold $k$, define the complement $I_k^c = \{1,\ldots,n\} \setminus I_k$.
	\item Using only data in $I_k^c$, estimate the nuisance functions $\hat{g}_k$ and $\hat{m}_k$.
	\item For observations in fold $k$, compute the orthogonalized scores.
	\item Average across all folds to obtain the final estimate.
\end{enumerate}

\subsection{Neyman Orthogonality Condition}

The score function $\psi(W; \theta, \eta)$ satisfies Neyman orthogonality if:

\[
\partial_\eta \mathbb{E}[\psi(W; \theta_0, \eta_0)] [\eta - \eta_0] = 0
\]

For the partially linear model, the orthogonal score is:

\[
\psi(W; \theta, \eta) = (Y - \theta D - g(X))(D - m(X))
\]

This score satisfies the orthogonality condition, making the estimator robust to small errors in nuisance function estimation.

\subsection{Asymptotic Variance of the DML Estimator}

Under regularity conditions, the DML estimator $\hat{\theta}_{\text{DML}}$ satisfies:

\[
\sqrt{n} (\hat{\theta}_{\text{DML}} - \theta_0) \xrightarrow{d} N(0, V)
\]

where the asymptotic variance $V$ is given by:

\[
V = \frac{\mathbb{E}[\varepsilon^2 \eta^2]}{(\mathbb{E}[\eta^2])^2}
\]

In practice, we estimate the variance using:

\[
\hat{\sigma}^2 = \frac{1}{n} \sum_{i=1}^n (\tilde{Y}_i - \hat{\theta}_{\text{DML}} \tilde{D}_i)^2
\]
\[
\widehat{\text{SE}}(\hat{\theta}_{\text{DML}}) = \sqrt{\frac{\hat{\sigma}^2}{\frac{1}{n} \sum_{i=1}^n \tilde{D}_i^2}}
\]

where $\tilde{Y}_i = Y_i - \hat{g}(X_i)$ and $\tilde{D}_i = D_i - \hat{m}(X_i)$ are the orthogonalized scores.

\section{Algorithm Pseudocode}

\subsection{DML Algorithm}

\begin{verbatim}
	Algorithm: Double Machine Learning (DML) for Partially Linear Model
	
	Input: Data {Y_i, D_i, X_i}_{i=1}^n, number of folds K
	Output: Estimate of treatment effect theta_hat
	
	1. Randomly partition data into K folds of equal size
	
	2. For k = 1 to K:
	a. Let I_k be indices in fold k, I_k^c be indices in complement
	b. Using only I_k^c, estimate:
	- g_hat_k(X) = E[Y | X] using random forest
	- m_hat_k(X) = E[D | X] using random forest
	c. For i in I_k, compute:
	- Y_tilde_i = Y_i - g_hat_k(X_i)
	- D_tilde_i = D_i - m_hat_k(X_i)
	
	3. Compute theta_hat = (sum_{i=1}^n D_tilde_i^2)^{-1} sum_{i=1}^n D_tilde_i Y_tilde_i
	
	4. Compute standard errors:
	sigma_hat_sq = (1/n) sum_{i=1}^n (Y_tilde_i - theta_hat D_tilde_i)^2
	SE(theta_hat) = sqrt(sigma_hat_sq / ((1/n) sum_{i=1}^n D_tilde_i^2))
	
	Return: theta_hat, SE(theta_hat), and p-value = 2(1 - Phi(|theta_hat|/SE(theta_hat)))
\end{verbatim}

\section{Summary of All Figures and Tables}

Table~\ref{tab:appendix_figures} provides a complete list of all figures generated in the analysis.

\begin{table}[htbp]
	\centering
	\caption{Complete List of Generated Figures}
	\label{tab:appendix_figures}
	\begin{tabular}{@{}ll@{}}
		\toprule
		\textbf{Figure} & \textbf{Description} \\
		\midrule
		dag\_causal\_diagram\_base.png & Base R DAG visualization \\
		dag\_causal\_diagram\_ggraph.png & ggraph DAG visualization \\
		psm\_glmm\_love\_plot\_*.png & Love plots for GLMM-based PSM balance assessment \\
		iptw\_glmm\_weights\_*.png & IPTW stabilized weight distributions \\
		dr\_glmm\_potential\_outcomes\_*.png & Doubly robust potential outcomes plots \\
		ps\_overlap\_glmm\_*.png & Propensity score overlap plots \\
		glmm\_caterpillar\_*.png & GLMM random effects caterpillar plots \\
		week4\_odds\_ratios\_comparison.png & Forest plot comparing traditional methods \\
		week5\_doubleml\_comparison.png & Forest plot comparing DML with traditional methods \\
		week6\_heterogeneity\_forest.png & Heterogeneity forest plot by subgroup \\
		week6\_robustness\_comparison.png & Robustness comparison across learners \\
		malaria\_by\_wealth.png & Malaria prevalence by wealth quintile \\
		itn\_by\_wealth.png & ITN use by wealth quintile \\
		malaria\_by\_residence.png & Malaria prevalence by urban/rural \\
		malaria\_by\_age.png & Malaria prevalence by age group \\
		itn\_by\_age.png & ITN use by age group \\
		malaria\_map\_2015.png & Geographic map of malaria prevalence (2015) \\
		malaria\_map\_2020.png & Geographic map of malaria prevalence (2020) \\
		itn\_map\_2015.png & Geographic map of ITN use (2015) \\
		itn\_map\_2020.png & Geographic map of ITN use (2020) \\
		correlation\_matrix.png & Correlation matrix of key variables \\
		\bottomrule
	\end{tabular}
\end{table}

\section{Computational Environment}

The analysis was conducted using R version 4.2.2. Key packages used include:

\begin{verbatim}
	- dplyr (1.1.0)      : Data manipulation
	- tidyr (1.3.0)      : Data cleaning
	- ggplot2 (3.4.1)    : Visualizations
	- survey (4.2)       : Survey-weighted analysis
	- lme4 (1.1-31)      : GLMM estimation
	- MatchIt (4.5.1)    : Propensity score matching
	- cobalt (4.4.1)     : Balance assessment
	- DoubleML (0.5.0)   : Double Machine Learning
	- mlr3 (0.15.0)      : Machine learning framework
	- ranger (0.14.1)    : Random forest
	- xgboost (1.7.3.1)  : Gradient boosting
\end{verbatim}

All code used to generate the results in this thesis is available from the author upon reasonable request. The analysis was performed on a standard laptop computer with 16 GB RAM and took approximately 2 hours to run all models and generate all figures.

\section{Declaration of AI Usage}

\subsection{Statement on Generative AI}

In accordance with AIMS academic integrity policies, the following statement is provided regarding the use of generative artificial intelligence (AI) tools in the preparation of this essay:

\begin{itemize}
	\item \textbf{Code Generation}: Generative AI (specifically, DeepSeek) was used to assist in writing and debugging R code for data cleaning, analysis, and visualization. All generated code was reviewed, tested, and validated by the author to ensure correctness and reproducibility.
	
	\item \textbf{Writing Assistance}: AI tools were used to assist with grammar checking, phrasing suggestions, and LaTeX formatting. All scientific content, interpretations, and conclusions remain the original work of the author.
	
	\item \textbf{No Plagiarism}: AI was not used to generate text that directly reproduces existing published work without proper citation.
	
	\item \textbf{Supervisor Approval}: The use of AI tools was disclosed to and approved by the supervisor, Dr. Haisa Osmanli.
\end{itemize}

The author takes full responsibility for the accuracy, integrity, and originality of the work presented herein.